\section{Results}
\label{sec:results}

\subsection{Main Results}

\begin{table}
  \caption{Main results on the validation set.}
  \label{tab:results}
  \begin{autotable}{L R R R}
    \toprule
    \textbf{Model} & \textbf{Metric 1} & \textbf{Metric 2} & \textbf{Metric 3} \\
    \midrule
    Baseline & 0.000 & 0.000 & 0/0/0 \\ \rowline
    Model A  & 0.000 & 0.000 & 0/0/0 \\ \rowline
    \textbf{Final} & \textbf{0.000} & \textbf{0.000} & \textbf{0/0/0} \\
    \bottomrule
  \end{autotable}
\end{table}

State the main finding of \cref{tab:results} first, then the details.

\subsection{Ideas Tried and Not Adopted}

% Text-heavy table: fixed-width wrapping columns
\begin{table}
  \caption{Ideas tried and not adopted.}
  \label{tab:not-adopted}
  \begin{autotable}{L L L}
    \toprule
    \textbf{Idea} & \textbf{Result} & \textbf{Decision} \\
    \midrule
    Idea one, described briefly & What happened on val & Kept / dropped \\ \rowline
    Idea two & What happened on val & Kept / dropped \\
    \bottomrule
  \end{autotable}
\end{table}

\subsection{Success and Failure Cases}

\begin{figure}
  \includegraphics[width=0.6\linewidth]{example-image-a}
  \caption{A success case (60\% width). Explain the colours and numbers in the
           image; a longer caption like this runs the full text width.}
  \label{fig:success}
\end{figure}
